\documentclass[11pt,a4paper]{article}
\usepackage[T1]{fontenc}
\usepackage[utf8]{inputenc}
\usepackage{geometry}
\usepackage{parskip}
\usepackage{hyperref}
\geometry{top=25mm, bottom=25mm, left=25mm, right=25mm}
\hypersetup{colorlinks=true, urlcolor=blue, linkcolor=black}
\begin{document}

Maryam Jalili \\
Department of Mathematics \\
Neyshabur Branch, Islamic Azad University \\
Neyshabur, Iran \\
\href{mailto:mrmjlili@iau.ac.ir}{mrmjlili@iau.ac.ir}

12 October 2026

Editor-in-Chief \\
\emph{Journal of Applied Mathematics and Computing} \\
Springer

Dear Editor,

I am pleased to submit my manuscript, \textbf{``A Diagnostic Framework for Detecting False Fractionalization in Physics-Informed Neural Networks: A Three-Way Comparison of Fractional Advection--Dispersion Residuals,'} for consideration for publication in the \emph{Journal of Applied Mathematics and Computing}.

\textbf{Motivation and contribution.} Physics-informed neural networks (PINNs) are increasingly deployed for fractional partial differential equations, yet the field lacks a systematic method to verify that a chosen residual discretization genuinely preserves the non-local memory of the underlying fractional operator. This paper addresses that gap through three contributions:

\begin{enumerate}
\item A \emph{three-test diagnostic protocol}---kernel-presence, $\Delta t$-refinement, and effective-time---that classifies a PINN residual as either a genuine fractional discretization or a local surrogate.
\item A \emph{local-limit lemma} for the naive two-point finite-difference approximation of the Atangana--Baleanu--Caputo derivative, with an explicit first-order error bound and a discussion of its scope (smooth test fields) relative to the singular step-inlet data used in the experiments.
\item The first \emph{three-way comparison} of PINN residuals guided by an explicit diagnostic protocol: naive FD (local), Caputo-L1 (power-law non-local), and ABC-L1 (Mittag--Leffler non-local), each trained under identical architecture, optimizer, epoch budget, and random seeds. The study reports MSE, training time, and the local-limit diagnostic at two fractional orders ($\alpha = 0.8$ and $\alpha = 0.5$), with careful discussion of the limits of the cross-PDE comparison.
\end{enumerate}

\textbf{Reproducibility.} All code, reference solutions, PINN checkpoints, and figure-generation scripts are released as a reproducible benchmark (GitHub: \url{https://github.com/mrmlili/aria-icnabs}; Zenodo DOI: \href{https://doi.org/10.5281/zenodo.23266644}{10.5281/zenodo.23266644}).

\textbf{Fit to JAMC.} The manuscript is a computational and applied mathematics contribution that combines numerical analysis of fractional discretizations with a machine-learning solver for time-fractional PDEs. The diagnostic protocol addresses a structural implementation issue (false fractionalization) that has not been systematically treated in the PINN literature. The work falls squarely within the scope of JAMC: applied numerical analysis of fractional operators, computational schemes for time-fractional evolution equations, and reproducible benchmarks for neural-network solvers of PDEs.

\textbf{Honest limitations.} The manuscript openly acknowledges three limitations and reports them in Section 7.2: (i) no manufactured-solution verification of the reference solvers; (ii) no formal consistency/convergence analysis of the ABC-L1 quadrature; and (iii) small seed ensembles (n=3). These limitations do not undermine the central contributions---the diagnostic protocol and the local-limit lemma---but they bound the scope of the numerical claims. We have made every effort to state each numerical result with the appropriate caveats.

The manuscript is original, has not been published elsewhere, and is not under consideration by any other journal.

Thank you for your time and consideration.

Sincerely, \\
Maryam Jalili

\end{document}
